\begin{table}[ht]
\caption{\textbf{\biology example.} The positive document explains the concept of meristem, which is the mechanism by which plant cells develop into other tissues and organs. This provides critical support for why a tree chunk sprouts and grows after being cut.}
\centering
% \resizebox{\textwidth}{!}{
\begin{tabular}{p{13cm}}
\toprule
% \hspace{5.5cm} Biology \\
% \midrule
\textit{\textbf{Query}} \\
\midrule
After I cut trees into logs and remove the branches in winter, they start growing. They sprout out and grow completely normal looking stems and leaves and maintain them all summer. The sprouts mostly appear around the cut branches. Sometimes they last all winter and grow for another year. How does a tree trunk sprout and grow after being cut?\\
\midrule
\textit{\textbf{Chain-of-thought reasoning to find documents}} \\
\midrule
The essential question is to figure out the growing mechanism of trees after being cut. It is usually related to cell division, which later develops into tissues and organs. We need to find relevant information of these cells in plants and their growing capabilities after being cut. \\
\midrule
\textit{\textbf{Example positive document}} \\
\midrule
Meristem \\
Tunica-corpus model of the apical meristem (growing tip). The epidermal (L1) and subepidermal (L2) layers form the outer layers called the tunica. The inner L3 layer is called the corpus. Cells in the L1 and L2 layers divide in a sideways fashion, which keeps these layers distinct, whereas the L3 layer divides in a more random fashion. \\
In cell biology, the meristem is a type of tissue found in plants. It consists of undifferentiated cells (meristematic cells) capable of cell division. Cells in the meristem can develop into all the other tissues and organs that occur in plants. These cells continue to divide until they become differentiated and lose the ability to divide. \\
There are three types of meristematic tissues: apical (at the tips), intercalary or basal (in the middle), and lateral (at the sides also known as cambium). At the meristem summit, there is a small group of slowly dividing cells, which is commonly called the central zone. Cells of this zone have a stem cell function and are essential for meristem maintenance. The proliferation and growth rates at the meristem summit usually differ considerably from those at the periphery. \\
... \\
Under appropriate conditions, each shoot meristem can develop into a complete, new plant or clone. Such new plants can be grown from shoot cuttings that contain an apical meristem. \\
... \\
\midrule
\textit{\textbf{Example negative document}} \\
\midrule
Cutting firewood from felled trees safely and easily (bucking) \\
There are loads of articles on the web showing how (and sometimes how not to!) cut down trees, but not so much about actually cutting firewood after the tree has been felled and stripped of its branches (called snedding or limbing in arborist speak). This ‘how to’ shows the easiest way to produce firewood by cutting many logs at the same time. I actually use my firewood cutting operation as part of my fitness regime, much more useful than going to the gym and the same muscle burn the day after! \\
I start by cutting the trees into 10’ (3m) logs so that I can easily drag them to my cleared working area close to the pick up point. \\
You can make a homemade ‘cradle’ that holds a whole bunch of logs from sturdy construction lumber (2 by 4). This enables you to safely cut many at once, as long as you put some of the heavy logs both at the bottom and on the top. Don’t forget to make the cradle narrow enough for the size of chainsaw bar you have. Usually not more than 12” or 300mm inside measurement. Make sure that it is constructed in such a way that you cannot hit any nails etc with the chainsaw. \\
Place the logs onto the cradle keeping the ends off the logs flush one side ( or at least staggered in 12” steps). If you are fussy like me you can mark one of the top logs at 12” (300mm) intervals with chalk to ensure even log lengths. Then cut alternate sides (so that it doesn’t topple over to one side), finishing off with a couple of cuts over the cradle itself, don’t forget to stop before you saw the thing in half lol! (See upgraded version below…) \\
... \\
\bottomrule
\end{tabular}
% }
\label{tab:biology_example}
\end{table}